% Part: proof-theory
% Chapter: natural-deduction
% Section: quantifiers

\documentclass[../../../include/open-logic-section]{subfiles}

\begin{document}

\olfileid{pt}{nat}{qua}

\olsection{Regular \usetoken{P}{proof} and Substitution}

The rules for the quantifiers introduce some complications owing to
the ``eigenvariable conditions'' that apply to the \Elim{\lexists}
and~\Intro{\lforall} rules. Most of these complications can be dealt
with by ensuring that the !!{constant}~$c$ in these inferences is
never reused. Let's introduce the following definition.

\begin{defn}
!!^a{proof}~$\delta$ in \Log{N1ci} or \Log{N2ci} is \emph{regular} if
\begin{enumerate}
  \item no eigenvariable~$a$ is used as the eigenvariable of more than
  one \Elim{\lexists} or~\Intro{\lforall} inference, and
  \item if $a$ is the eigenvariable of an \Elim{\lexists}
  or~\Intro{\lforall} inference, it only occurs in the immediate
  sub-!!{proof} leading to the minor premise, or the premise, respectively.
\end{enumerate}
\end{defn}

\begin{lem}\ollabel{lem:subst-regular} If $\delta$ is regular, ends
in~$!C$, $c$ is !!a{constant} not used as an eigenvariable
in~$\delta$, and $t$~is a term not containing any eigenvariables
of~$\delta$, then $\Subst{\delta}{t}{c}$ resulting from~$\delta$ by
replacing every occurrence of~$c$ by~$t$ is a regular !!{proof}. The
end-!!{formula} of $\Subst{\delta}{t}{c}$ is $\Subst{!C}{t}{c}$. If
$!A^x$ is an open assumption of~$\delta$ then $\Subst{!A}{t}{c}^x$ is
an open assumption of~$\Subst{\delta}{t}{c}$.
\end{lem}

\begin{proof}
  By induction on $\pheight{\delta}$. If $\pheight{\delta} = 0$, it
  consists simply of an assumption~$!A^x$. Then $\Subst{\delta}{t}{c}$
  is~$\Subst{!A}{t}{c}^x$.

  If $\pheight{\delta} > 0$ we distinguish cases according to the last
  inference in~$\delta$. The cases for rules for propositional
  connectives are easy and left as exercises. We consider the cases
  where $\delta$ ends in a quantifier inference.
  \begin{enumerate}
    \item If the last inference is $\Intro{\lforall}$, $\delta$ is of
    the form
    \[
    \AxiomC{}
    \RightLabel{$\delta'$}
    \DeduceC{$!A(a,c)$}
    \RightLabel{\Intro{\lforall}}
    \UnaryInfC{$\lforall[x][!A(x,c)]$}
    \DisplayProof
    \]
    By induction hypothesis, $\Subst{\delta'}{t}{c}$ is a regular
    !!{proof} of $!A(a,t)$. Since $a$ is an eigenvariable of~$\delta$,
    $a$~does not occur in~$t$. Therefore, the last inference in
    $\Subst{\delta}{t}{c}$,
    \[
    \AxiomC{}
    \RightLabel{$\Subst{\delta'}{t}{c}$}
    \DeduceC{$!A(a,t)$}
    \RightLabel{\Intro{\lforall}}
    \UnaryInfC{$\lforall[x][!A(x,t)]$}
    \DisplayProof
    \]
    is correct: since $t$ does not contain~$a$, neither the conclusion
    nor any open assumption contains~$a$.
    \item If the last inference is $\Elim{\lforall}$,
    $\delta$ is of the form
    \[
    \AxiomC{}
    \RightLabel{$\delta'$}
    \DeduceC{$\lforall[x][!A(x,c)]$}
    \RightLabel{\Elim{\lforall}}
    \UnaryInfC{$!A(s(c),c)$}
    \DisplayProof
    \]
    By induction hypothesis, $\Subst{\delta'}{t}{c}$ is a regular
    !!{proof} of $\lforall[x][!A(x,t)]$. (The term~$s(c)$ in the
    conclusion may, but need not, contain~$c$.) The last inference in
    $\Subst{\delta}{t}{c}$,
    \[
    \AxiomC{}
    \RightLabel{$\Subst{\delta'}{t}{c}$}
    \DeduceC{$\lforall[x][!A(x,t)]$}
    \RightLabel{\Elim{\lforall}}
    \UnaryInfC{$!A(s(t),t)$}
    \DisplayProof
    \]
    is correct, and $!A(s(t),t) \ident \Subst{!A(s(c),c)}{t}{c}$. The
    case for \Intro{\lexists} is similar.
  \item If the last inference is $\Elim{\lexists}$, $\delta$ is of
    the form
    \[
    \AxiomC{}
    \RightLabel{$\delta_1'$}
    \DeduceC{$\lexists[x][!A(x,c)]$}
    \AxiomC{$\Discharge{!A(a,c)}{x}$}
    \RightLabel{$\delta_2'$}
    \DeduceC{$!C$}
    \DischargeRule{\Elim{\lexists}}{x}
    \BinaryInfC{$!C$}
    \DisplayProof
    \]
    By induction hypothesis, $\Subst{\delta_1'}{t}{c}$ and
    $\Subst{\delta_2'}{t}{c}$ are regular !!{proof}s of
    $\lexists[x][!A(x,c)]$ and of $!C$ from the open
    assumption~$!A(a,t)$. Since $a$ is an eigenvariable of~$\delta$,
    $a$~does not occur in~$t$. Therefore, the last inference in
    $\Subst{\delta}{t}{c}$,
    \[
    \AxiomC{}
    \RightLabel{$\Subst{\delta'}{t}{c}$}
    \DeduceC{$!A(a,t)$}
    \RightLabel{\RightR{\lforall}}
    \UnaryInfC{$\lforall[x][!A(x,t)]$}
    \DisplayProof
    \]
    is correct: $\Subst{!A(x,t)}{a}{x} \ident \Subst{!A(a,c)}{t}{c}$, neither the conclusion~$!C$ nor any open assumption contains~$a$. 
  \end{enumerate}
\end{proof}

\begin{prop}\ollabel{prop:regular}
If $\delta$ is !!a{proof} in \Log{N1c} or \Log{N1i}, then there is a
regular !!{proof}~$\delta'$ with the same structure (in particular, the
same !!{height}) which differs from $\delta$ only by replacing
eigenvariables.
\end{prop}

\begin{proof}
For the purpose of the proof only, call an eigenvariable inference
of~$\delta$ \emph{clean} if its eigenvariable is not the eigenvariable of
another inference nor does it occur anywhere in~$\delta$ except the
immediate sub-!!{proof} of the inference, and \emph{dirty} otherwise.
Let $n$ be the number of dirty eigenvariable inferences in~$\delta$. We
show, by induction on~$n$, that $\delta$~can be transformed into the
required regular !!{proof}~$\delta'$. If $n=0$, all eigenvariable
inferences in~$\delta$ are clean, and so $\delta$ is regular, and there is
nothing to prove. Otherwise, pick a highest eigenvariable inference,
say \Intro{\lforall}. The sub-!!{proof}~$\delta_1$ of~$\delta$ ending in it
is of the form
    \[
    \AxiomC{}
    \RightLabel{$\delta_2$}
    \DeduceC{$!A(c)$}
    \RightLabel{\Intro{\lforall}}
    \UnaryInfC{$\lforall[x][!A(x)]$}
    \DisplayProof
    \]
Note that as $c$ is an eigenvariable, it does not occur
in~$\lforall[x][!A(x)]$ or an open assumption. The proof $\delta_2$ is
regular, since it contains no eigenvariable inferences at all. Let $a$
be !!a{constant} not occuring in~$\delta$. By
\olref{lem:subst-regular}, $\Subst{\delta_2}{a}{c}$ is a regular proof
of $!A(a)$. By the eigenvariable condition, no open assumption
of~$\delta_2$ contains~$c$, hence no open assumption
of~$\Subst{\delta_2}{a}{c}$ contains~$a$. Thuse we can
apply~$\Intro{\lforall}$. If we replace $\delta_1$ in~$\delta$ by the
proof $\delta_1'$,
    \[
    \AxiomC{}
    \RightLabel{$\Subst{\delta_2}{a}{c}$}
    \DeduceC{$!A(a)$}
    \RightLabel{\Intro{\lforall}}
    \UnaryInfC{$\lforall[x][!A(x)]$}
    \DisplayProof
    \]
we have replaced a dirty eigenvariable inference by a clean one, and
the only difference is the replacement of the eigenvariable $c$
by~$a$. The resulting proof has $n-1$ dirty eigenvariable inferences,
and the result follows from the inductive hypothesis.
\end{proof}

\begin{prob}
Describe the case in the proof of \olref{prop:regular} where the
topmost eigenvariable inference is~\Elim{\lexists}.
\end{prob}

We will not appeal to the proposition just proved explicitly, but we
will assume that all !!{proof}s considered are regular---if they are
not, we can always replace eigenvariables to make them regular.

\Olref{lem:subst-regular} and \olref{prop:regular} also hold for
\Log{N2ci}. We leave the proofs as exercises.

\end{document}